\chapter{Дискретность микроуровня}
\label{ch:foundations}

\section{От постулата к теореме}

Ранний постулат «решётка существует'' снят.
При абсолютных условиях (гл.~\ref{ch:axioms}) и геометрии носителя (гл.~\ref{ch:carrier})
дискретность~$\carrierLattice$, шаг~$\caStepSpace$ и такт~$\htick$ \emph{выводятся}.

\begin{lemma}[Счётное пространство конфигураций]
\label{lem:countable-config}
В постановке поля и нормы (\cref{def:field-a3}) пространство микросостояний
\[
  \mathcal{C}=\bigl\{z(x)\in\mathbb{C}^2\bigr\}_{x\in\carrierLattice}
\]
с дискретным алфавитом на каждом узле счётно, а эволюция $g:\mathcal{C}\to\mathcal{C}$ --- биекция.
\end{lemma}

\begin{proof}
\pusing{определения поля (\cref{def:field-a3}) и бюджета (\cref{sec:bit-budget})}
на узле~$x$ состояние лежит в конечном кольце $\mathbb{Z}_{N_{\mathrm{ring}}}[i]$;
мощность кольца фиксирует $B_{\hv}$.
\phence при счётном~$\carrierLattice$ пространство
\[
  \mathcal{C}=\bigl\{z(x)\in\mathbb{C}^2\bigr\}_{x\in\carrierLattice}
\]
есть счётное произведение конечных множеств и само счётно.
\pfrom{сохранения нормы и обратимости шага (\Axref{ax:A13})}
$g$ --- перестановка (биекция) на~$\mathcal{C}$.
\end{proof}

\begin{lemma}[\latGCKshort: носитель и элементарная шкала]
\label{lem:fcc-scale}
Носитель --- \latGCK\ $\carrierLattice$ с $|N(x)|=12$ (\cref{thm:fcc-choice,def:packing-types}).
Шаг~$\caStepSpace$ и такт~$\htick$ однозначно задаются однотактовым телом на \latGCKshort\ (\cref{sec:fcc-kappa});
$\hv$ --- узел решётки (регистр поля и нормы (\cref{def:field-a3})).
\end{lemma}

\begin{proof}
\pfrom{теоремы~\ref{thm:fcc-choice}}
\[
  \carrierLattice\ \text{--- \latGCK},\qquad |N(x)|=12.
\]

\smallskip
\emph{Шаг~$\caStepSpace$.}
\pusing{геометрии однотактового тела (\cref{sec:fcc-kappa})}
пусть $a=\caStepSpace$ --- ребро кубооктаэдра.
\pfollows
объём ячейки Вороного
\[
  v_{\hv}=\frac{a^3}{\sqrt{2}}=\frac{\caStepSpace^3}{\sqrt{2}},
\]
минимальный регистр при плотнейшей упаковке.

\smallskip
\emph{Шаг~$\htick$.}
Тот же раздел задаёт $\caKappaGeom=R_{\mathrm{in}}/R_{\mathrm{out}}=1/\sqrt{2}$ и
$\caSignalSpeed=\caStepSpace/\htick$.
\pfrom{планковской идентификации (\cref{def:si-bridge}) при $\ell_P:=\caStepSpace$}
\begin{align}
  \htick &= \caKappaGeom\,t_P, \\
  \caSignalSpeed &= \sqrt{2}\,\speedC.
\end{align}

\smallskip
\emph{Ячейка~$\hv$.}
\puseref{определения поля и нормы}{def:field-a3}
получаем: $\PlanckCell$ --- узел \latGCKshort\ (регистр~$z$); $dV$ --- объём Вороного ($V_P/\sqrt{2}$ при $a=\ell_P$).
\end{proof}

\begin{lemma}[Атомность регистра]
\label{lem:register-atom}
Элементарная ячейка~$\hv$ не дробится
на под-решётку с шагом $\caStepSpace'<\caStepSpace$ без смены алфавита и закона~$g$.
\end{lemma}

\begin{proof}
\pfrom{запрета двух совпадающих спиноров в одном~$\PlanckCell$}
регистр на узле атомарен.
\pusing{минимального фазового шага $\Delta\phi_{\min}=\tfrac12$}
\begin{align}
  \ladderAction &= \frac{\planckHbar}{2}, \\
  \ladderMomentum &= \frac{\planckHbar}{2\caStepSpace}.
\end{align}
Топологический заряд $n\in\mathbb{Z}$ и полюс привязаны к границе одной ячейки~$\hv$;
размазывание полюса по подузлам разнесло бы $n$.
Потолок плотности и $|z|\ge z_{\min}>0$ фиксируют бюджет~$B_{\hv}$ на регистр.
\phence формальное дробление~$\hv$ потребовало бы нескольких независимых алфавитов на месте одного узла \latGCKshort, то есть другого закона~$g$.
\end{proof}

\begin{theorem}[Дискретность и минимальная ячейка $\hv$]
\label{thm:discreteness}
Пусть $g$ --- допустимый локальный закон на~$\layerMicro$,
удовлетворяющий абсолютным условиям~\cref{ch:axioms} на геометрии \latGCKgen~\cref{ch:carrier}.
Тогда:
\begin{enumerate}[label=(\roman*)]
\item $g$ --- дискретная динамика на~$\carrierLattice$;
\item существуют счётный носитель~\latGCK\ $\carrierLattice$, элементарная ячейка~$\hv$, шаг~$\caStepSpace$ и такт~$\htick$;
\item дробление регистра «ниже~$\caStepSpace$'' меняет алфавит: под-решётка несёт другой закон.
\end{enumerate}
\end{theorem}

\begin{proof}
\emph{(i).}
Предположим, $g$ --- «поле на континууме'': $z(x)$ при $x\in\mathbb{R}^3$ с несчётными степенями свободы в точке.
\pfrom{леммы~\ref{lem:countable-config}}
допустимый закон действует на счётном $\mathcal{C}$ с конечным алфавитом на узле.
Биективная эволюция на $\mathcal{C}$ не воспроизводит континуальный поток на~$\mathbb{R}^3$.
\phence $g$ --- дискретная динамика на~$\carrierLattice$.

\smallskip
\emph{(ii).}
\pfrom{теоремы~\ref{thm:fcc-choice} и леммы~\ref{lem:fcc-scale}}
существуют \latGCK\ $\carrierLattice$, $|N|=12$, шаг $\caStepSpace$, такт $\htick$ и ячейка~$\hv$.

\smallskip
\emph{(iii).}
Если подрешётка с $\caStepSpace'<\caStepSpace$ несла тот же~$g$, на узле \latGCKshort\ было бы несколько регистров с общим топологическим зарядом.
\pfrom{леммы~\ref{lem:register-atom}}
это запрещено: заряд и полюс привязаны к одной~$\hv$, бюджет не делится без смены алфавита.
\end{proof}

\begin{corollary}[Дискретность как следствие]
\label{cor:discreteness}
Дискретность микроуровня есть содержание теоремы~\ref{thm:discreteness}.
\end{corollary}

\section{Состояние, континуум, вакуум}

\begin{corollary}[Континуум на~$M$ отсутствует]
\label{cor:no-continuum}
Непрерывные поля и предел $\caStepSpace\to 0$ --- свойство макроуровня~$T$
(гл.~\ref{sec:no-continuum-carrier}).
\end{corollary}

\begin{corollary}[Нижняя граница амплитуды]
\label{cor:vacuum-floor}
$|z(x)|\ge z_{\min}>0$ на~$M$.
\end{corollary}

\section{Замыкание закона, механика, кванты}

\begin{corollary}[Замыкание $g$]
\label{cor:g-closure}
Допустимый закон~$g$ --- закон из теоремы о законе $g$ (\cref{thm:g-law}).
\end{corollary}

\begin{corollary}[Лестница Ландау]
\label{cor:landau}
представление Маделунга --- язык~$T$.
\end{corollary}

\begin{corollary}[Дискретные возбуждения]
\label{cor:full-quantization}
На~$M$ --- занятость мод $n_k\in\mathbb{Z}$; «волны'' --- макроскопическое осреднение на~$T$.
\end{corollary}


